\chapter{Source Code Listings}
\label{app:code}

This appendix contains the main source files of the Gift Advisor Bot. The full project (including experiment scripts and 30 personas) is available on GitHub.

\section{Database Module (\texttt{database.py})}

\begin{lstlisting}[language=Python, caption={SQLite database module.}, label={lst:app-db}]
import json
import sqlite3
from pathlib import Path

DB_PATH = Path("gift_bot.db")


def _conn() -> sqlite3.Connection:
    conn = sqlite3.connect(DB_PATH)
    conn.row_factory = sqlite3.Row
    return conn


def init_db() -> None:
    with _conn() as c:
        c.executescript("""
            CREATE TABLE IF NOT EXISTS users (
                id INTEGER PRIMARY KEY AUTOINCREMENT,
                telegram_id INTEGER NOT NULL UNIQUE,
                username TEXT,
                registered_at TEXT DEFAULT CURRENT_TIMESTAMP
            );
            CREATE TABLE IF NOT EXISTS sessions (
                id INTEGER PRIMARY KEY AUTOINCREMENT,
                user_id INTEGER NOT NULL,
                profile_data_json TEXT NOT NULL,
                strategy_used TEXT NOT NULL,
                response_json TEXT,
                timestamp TEXT DEFAULT CURRENT_TIMESTAMP,
                FOREIGN KEY (user_id) REFERENCES users(id)
            );
            CREATE TABLE IF NOT EXISTS feedback (
                id INTEGER PRIMARY KEY AUTOINCREMENT,
                session_id INTEGER NOT NULL,
                rating_relevance INTEGER,
                rating_creativity INTEGER,
                rating_specificity INTEGER,
                comment TEXT,
                FOREIGN KEY (session_id) REFERENCES sessions(id)
            );
        """)


def save_user(telegram_id: int, username: str) -> int:
    with _conn() as c:
        c.execute("INSERT OR IGNORE INTO users (telegram_id, username) "
                  "VALUES (?, ?)", (telegram_id, username))
        row = c.execute("SELECT id FROM users WHERE telegram_id = ?",
                        (telegram_id,)).fetchone()
        return row["id"]


def save_session(telegram_id, profile, strategy, response):
    user_id = save_user(telegram_id, "")
    with _conn() as c:
        cur = c.execute("INSERT INTO sessions "
                        "(user_id, profile_data_json, strategy_used, response_json) "
                        "VALUES (?, ?, ?, ?)",
                        (user_id, json.dumps(profile), strategy, json.dumps(response)))
        return cur.lastrowid


def save_feedback(session_id, rating_relevance, rating_creativity,
                  rating_specificity, comment):
    with _conn() as c:
        c.execute("INSERT INTO feedback (session_id, rating_relevance, "
                  "rating_creativity, rating_specificity, comment) "
                  "VALUES (?, ?, ?, ?, ?)",
                  (session_id, rating_relevance, rating_creativity,
                   rating_specificity, comment))
\end{lstlisting}

\section{Grok API Client (\texttt{grok\_client.py})}

\begin{lstlisting}[language=Python, caption={Async client for the Grok API.}, label={lst:app-grok}]
import os
import httpx
from dotenv import load_dotenv

load_dotenv()

GROK_API_KEY = os.getenv("GROK_API_KEY", "")
GROK_MODEL = os.getenv("GROK_MODEL", "grok-2-latest")
GROK_URL = "https://api.x.ai/v1/chat/completions"
TIMEOUT = 60.0


async def call_grok(prompt: str, temperature: float = 0.7) -> str:
    if not GROK_API_KEY:
        raise RuntimeError("GROK_API_KEY is missing. Set it in .env file.")
    headers = {
        "Authorization": f"Bearer {GROK_API_KEY}",
        "Content-Type": "application/json",
    }
    payload = {
        "model": GROK_MODEL,
        "messages": [{"role": "user", "content": prompt}],
        "temperature": temperature,
    }
    async with httpx.AsyncClient(timeout=TIMEOUT) as client:
        response = await client.post(GROK_URL, json=payload, headers=headers)
        response.raise_for_status()
        data = response.json()
        return data["choices"][0]["message"]["content"].strip()
\end{lstlisting}

\section{Prompt Builder with Four Strategies (\texttt{prompts/builder.py})}

\begin{lstlisting}[language=Python, caption={Prompt builder for the four strategies.}, label={lst:app-prompts}]
STRATEGIES = ["naive", "constrained", "cot", "persona"]


def _profile_block(profile: dict) -> str:
    interests = ", ".join(profile.get("interests", [])) or "not specified"
    free_text = profile.get("free_text") or "none"
    return (f"- Relation: {profile['relation']}\n"
            f"- Gender: {profile['gender']}\n"
            f"- Age: {profile['age_group']}\n"
            f"- Occasion: {profile['occasion']}\n"
            f"- Budget: {profile['budget']}\n"
            f"- Interests: {interests}\n"
            f"- Extra notes: {free_text}")


def build_naive_prompt(profile):
    return f"Suggest a gift for this person:\n{_profile_block(profile)}"


def build_constrained_prompt(profile):
    return f"""Suggest exactly 3 gift ideas for this person.

Profile:
{_profile_block(profile)}

Rules:
- Stay within the budget.
- Do NOT suggest flowers, chocolates, or generic gift cards.
- Each idea must be a specific product type.
- For each idea: Name, Why it fits (1 sentence), Approximate price in USD.
"""


def build_cot_prompt(profile):
    return f"""You are helping someone choose a gift. Think step by step.

Profile:
{_profile_block(profile)}

Step 1: What does this person likely enjoy?
Step 2: What types of gifts match those interests AND the budget?
Step 3: Filter out cliches (flowers, chocolates, generic cards).
Step 4: Pick 3 specific ideas.

Finally, write the 3 ideas as: Name | Why it fits | Price in USD
"""


def build_persona_prompt(profile):
    return f"""You are Elena, an expert personal gift advisor with
20 years of experience. You hate cliches.

Gift receiver profile:
{_profile_block(profile)}

Give your 3 best recommendations. For each idea:
- Name of the gift
- One sentence why it fits this person
- Approximate price in USD
"""


BUILDERS = {
    "naive": build_naive_prompt,
    "constrained": build_constrained_prompt,
    "cot": build_cot_prompt,
    "persona": build_persona_prompt,
}


def build_prompt(strategy: str, profile: dict) -> str:
    if strategy not in BUILDERS:
        raise ValueError(f"Unknown strategy: {strategy}")
    return BUILDERS[strategy](profile)
\end{lstlisting}

\section{Bot Entry Point (\texttt{bot.py}, simplified)}

\begin{lstlisting}[language=Python, caption={Main bot entry point (key parts).}, label={lst:app-bot}]
import asyncio
import os
import time

from aiogram import Bot, Dispatcher, F
from aiogram.filters import CommandStart, Command
from aiogram.fsm.context import FSMContext
from aiogram.fsm.state import State, StatesGroup
from aiogram.fsm.storage.memory import MemoryStorage
from dotenv import load_dotenv

from database import init_db, save_session, save_feedback
from grok_client import call_grok
from prompts.builder import build_prompt

load_dotenv()
bot = Bot(token=os.getenv("BOT_TOKEN"))
dp = Dispatcher(storage=MemoryStorage())


class GiftForm(StatesGroup):
    relation = State()
    gender = State()
    age_group = State()
    occasion = State()
    budget = State()
    interests = State()
    free_text = State()


@dp.message(CommandStart())
async def cmd_start(message, state):
    await state.clear()
    await message.answer("Welcome to Gift Advisor Bot!")


async def generate_recommendation(message, state):
    data = await state.get_data()
    profile = {k: data[k] for k in ["relation", "gender", "age_group",
                                     "occasion", "budget", "interests"]}
    profile["free_text"] = data.get("free_text", "")

    strategy = os.getenv("DEFAULT_STRATEGY", "cot")
    prompt = build_prompt(strategy, profile)

    t0 = time.time()
    answer = await call_grok(prompt)
    elapsed = time.time() - t0

    save_session(message.from_user.id, profile, strategy,
                 {"answer": answer, "elapsed": elapsed})
    await message.answer(f"3 gift ideas:\n\n{answer}")


async def main():
    init_db()
    await dp.start_polling(bot)


if __name__ == "__main__":
    asyncio.run(main())
\end{lstlisting}

\section{Experiment Runner (\texttt{experiment/run\_experiment.py})}

\begin{lstlisting}[language=Python, caption={Experiment runner: 30 personas $\times$ 4 strategies.}, label={lst:app-exp}]
import asyncio
import json
import time
from pathlib import Path

from grok_client import call_grok
from prompts.builder import build_prompt, STRATEGIES

PERSONAS_PATH = Path("experiment/personas.json")
RESULTS_JSON = Path("experiment/results.json")
PAUSE = 1.0


async def run_one(persona, strategy):
    prompt = build_prompt(strategy, persona)
    t0 = time.time()
    try:
        answer = await call_grok(prompt)
        error = ""
    except Exception as exc:
        answer = ""
        error = str(exc)
    elapsed = time.time() - t0
    return {"persona_id": persona["id"], "strategy": strategy,
            "elapsed_seconds": round(elapsed, 2),
            "answer": answer, "error": error}


async def main():
    personas = json.loads(PERSONAS_PATH.read_text())
    results = []
    for persona in personas:
        for strategy in STRATEGIES:
            row = await run_one(persona, strategy)
            results.append(row)
            await asyncio.sleep(PAUSE)
            RESULTS_JSON.write_text(json.dumps(results, indent=2))
    print(f"Done. {len(results)} cases saved.")


if __name__ == "__main__":
    asyncio.run(main())
\end{lstlisting}
